\section{Introduction}
\label{sec:intro}

% ~1 page; final numbers filled by scripts/fill_results.py when results land.

Large language models routinely undergo sequential fine-tuning for alignment, domain adaptation, and chained task specialisation, yet each adaptation carries a known cost: \textit{catastrophic forgetting} \citep[e.g.,][]{mccloskey1989catastrophic, french1999catastrophic, kirkpatrick2017ewc}. Parameter-efficient fine-tuning (PEFT) methods such as LoRA \citep{hu2022lora} partially mitigate this phenomenon by constraining updates to a low-rank subspace, but they do so \textit{uniformly} across the transformer, adapting every queried, keyed, and MLP projection with the same rank and target set.

We argue that this uniformity is unwarranted. Mechanistic interpretability work has shown that transformer components serve strikingly different functional roles \citep{elhage2021mathematical, meng2022locating, geva2023dissecting}: early layers encode general linguistic features, some attention heads act as ``anchor'' circuits for named-entity resolution, and late-layer MLPs tend to store task-specific associations. If function varies by component, then \textit{vulnerability to forgetting should vary by component}.

This paper tests that hypothesis. Concurrent work has begun to analyse forgetting via gradient attribution \citep{anon2026mechforget, lrdistribution2024} and to localise attention heads for new-task adaptation \citep{yin2024lofit}. We complement this line in two directions: we use \textit{causal} intervention rather than gradient attribution, and we invert the localisation objective from adaptation to \textit{resistance}. Our contributions are:

\begin{itemize}
  \item The first \textbf{causally-traced} per-component forgetting topology for Qwen3-8B and Llama-3.1-8B across a four-task sequence, using ROME-style counterfactual intervention.
  \item The \textbf{Forgetting Vulnerability Score} (\fvs{}), a transferable metric that predicts component vulnerability before running downstream tasks.
  \item \textbf{\tglora{}}, the first PEFT method to operationalise a forgetting map: LoRA restricted to the $50\%$ most resistant components. \tglora{} matches standard LoRA on current tasks while reducing prior-task forgetting by \pend{forget_reduction_pct}.
  \item Empirical arbitration between gradient-attribution and causal-intervention accounts of \textit{which} components drive forgetting.
\end{itemize}

Our findings align with the EMNLP 2026 special theme on explainability: the forgetting topology \textit{is} an explanation. It answers not only \textit{whether} a fine-tune will erode prior capabilities, but \textit{where} in the network that erosion will occur.
